\chapter{Att kontrollera ett påstående}
\label{app:verifiera-villkor}

Det här materialet gör en rad påståenden.  De flesta går att pröva vid
datorn: körs båda grenarna eller inte?  Slutar programmet vid ett falskt
villkor eller inte?  Sådana påståenden ska du inte tro på, du ska prova
dem --- och materialet är byggt så att du kan: varje exempel är ett
program som körs, och utskriften under det är den riktiga.

Men några av påståendena handlar om något annat --- om vad
\emph{forskningen} säger om vad nybörjare brukar missförstå.  De går inte
att pröva vid datorn:
\begin{enumerate}
  \item att nybörjare skriver ett likhetstecken där en jämförelse avses,
    och läser en tilldelning som ett påstående om likhet
    (\cref{sec:villkor});
  \item att nybörjare tror att ett falskt villkor avslutar programmet,
    och att båda grenarna körs (\cref{sec:forgrening-villkor});
  \item att nybörjare kedjar villkor som
    \mintinline{python}{if x != a or b} (\cref{sec:villkor});
  \item att en slinga vars styrvariabel aldrig ändras är ett
    återkommande fel, att slingvariabeln glöms bort före slingan, och att
    en upprepning tas till där ett val räcker (\cref{sec:slingor})%
  \ifbool{ltnote}{;
  \item att en kontrast måste upplevas \emph{samtidigt} för att den
    varierande aspekten ska kunna urskiljas --- skälet till att exemplen
    står parvis (hela materialet).}{.}
\end{enumerate}
Hur vet man om sådant stämmer?  Samma fråga möter dig varje gång du läser
en lärobok, en webbsida eller ett svar från en språkmodell --- och varje
gång du själv ska skriva en rapport.  Arbetsgången är densamma i alla
kursens föreläsningar och beskrivs i föreläsningen \emph{Algoritmiskt
tänkande}, bilaga~A: hur ett påstående görs till en fråga som kan få
svaret nej, hur man väljer rätt sorts källa och söker ämnesdrivet och
efter invändningar, hur källan läses i fulltext och hur läsningen
antecknas, och hur svaret skrivs ut med sina förbehåll.  Här står bara
det som är särskilt för den här föreläsningen.

\Cref{app:missuppfattningar-villkor} tillämpar en del av arbetsgången på
de \ifbool{ltnote}{fem}{fyra} påståendena i listan: källorna är
återanvända ur en annan genomgång och lästa om, men någon ny sökning har
inte gjorts (\cref{sec:cond-metod}).  En återanvänd källa är inte redan
kontrollerad bara för att någon annan har citerat den; därför är varje
citat läst om i originalet.  Bilagan slutar med ett avsnitt
\emph{Fortsatt arbete}: vad som lämnades ogjort, och vad som skulle
behöva undersökas för ett säkrare svar.  \Cref{tab:paastaaenden-cond}
säger var svaret på varje påstående står, också på dem som är belagda i
andra föreläsningar eller ännu inte prövade.

\begin{table}[htbp]
  \centering
  \caption{Påståendena i den här föreläsningen som frågor, och var
    svaret står.}
  \label{tab:paastaaenden-cond}
  \begin{tabular}{@{}p{0.64\linewidth}p{0.32\linewidth}@{}}
    \toprule
    Fråga & Var svaret står \\
    \midrule
    Skriver nybörjare ett likhetstecken där en jämförelse avses, och
      läser de en tilldelning som ett likhetstest? &
      \Cref{app:missuppfattningar-villkor} \\
    Tror nybörjare att ett falskt villkor avslutar programmet, eller att
      båda grenarna körs? & \Cref{app:missuppfattningar-villkor} \\
    Kedjar nybörjare villkor utan att upprepa jämförelsen? &
      \Cref{app:missuppfattningar-villkor} \\
    Är oändliga slingor, oinitierade slingvariabler och upprepning
      i stället för val dokumenterade nybörjarfel? &
      \Cref{app:missuppfattningar-villkor} \\
    % Rows for claims made only in the instructor notes print only in
    % the teacher edition.
    \ifbool{ltnote}{%
    Måste en kontrast upplevas samtidigt för att kunna urskiljas? &
      \Cref{app:missuppfattningar-villkor} \\
    }{}%
    Är det bra att inte upprepa sig i ett program, och kommer principen
      \foreignlanguage{english}{DRY} från Hunt och Thomas? &
      Redan belagt i \emph{Algoritmiskt tänkande}, bilaga E \\
    Kommer principen om ett enda ansvar
      (\foreignlanguage{english}{single responsibility}) från Robert
      C. Martin, och är den etablerad praxis? & Redan belagt i
      \emph{Funktioner}, bilaga B \\
    \ifbool{ltnote}{%
    Uppfattar studenten en rekommendation som frivillig om materialet
      inte själv följer den? & Inte prövat; se \cref{sec:cond-fortsatt}
      \\
    }{}%
    \bottomrule
  \end{tabular}
\end{table}

\chapter{Är missuppfattningarna belagda?}
\label{app:missuppfattningar-villkor}
\chapterprecis{Författaren har ännu inte granskat resultaten i den här
  bilagan i sin helhet.}

Den här föreläsningen är byggd kring ett antal saker som nybörjare
brukar missförstå, och kontrastparen är valda för att göra just dem
synliga: samma program med ett sant och ett falskt villkor, samma
villkor med ett och två likhetstecken, samma uppgift med
\mintinline{python}{if} och med \mintinline{python}{while}.  Frågan den
här bilagan besvarar är: \emph{är de missuppfattningar materialet är
byggt kring faktiskt belagda, och räcker beläggen för det materialet
påstår?}

\section{Metod}
\label{sec:cond-metod}

Ingen ny litteratursökning har gjorts för den här föreläsningen.
Källorna är hämtade ur en systematisk genomgång av forskningen om
nybörjares missuppfattningar som kursens lärare gjort inom ett pågående
forskningsarbete\autocite{SooriBosk2026}.  Den genomgången söktes
ämnesdrivet i flera databaser
--- Semantic Scholar, OpenAlex, DBLP, arXiv, Web of Science, IEEE Xplore
och Scopus --- på frågan vilka missuppfattningar nybörjare har i
grundläggande programmering, och träffarna sorterades först maskinellt
och sedan för hand.

Att återanvända en källa är inte detsamma som att ha kontrollerat den.
För varje källa som citeras här har därför tre saker gjorts om:
påståendet har formulerats så som \emph{den här} föreläsningen gör det;
det ordagranna citat föreläsningen stöder sig på har lästs upp igen i
källan; och det har antecknats vilken nivå kontrollen nådde --- fulltext
eller enbart sammanfattning.  Anteckningarna står i materialets källfil
för litteraturlistan, ovanför varje post.

Den här bilagan innehåller därför ingen träfflista och inga sökfrågor:
det vore att låtsas att en sökning gjorts som inte gjorts.  Vad som står
här är i stället vilka källor som bär vilket påstående, och hur långt
kontrollen av dem räcker.

\section{Resultat}
\label{sec:cond-resultat}

\paragraph{Ett likhetstecken eller två.}  Att nybörjare skriver
\mintinline{python}{=} där \mintinline{python}{==} avses är dokumenterat
i en katalog över antimönster, byggd ur 95 missuppfattningar i verklig
kod från grundkursstudenter.  I katalogen står att
\textcquote{Bosse2021Antipatterns}{using only an equals sign '=' is
attribution and not a comparison} --- ordet
\enquote{\foreignlanguage{english}{attribution}} är källans egen
återgivning av portugisiskans \emph{atribuição}, alltså tilldelning.
Felet har en spegelbild i läsriktningen: i en intervjustudie av 41
deltagare läste studenterna en tilldelning som ett påstående om likhet
som avgjorde om programmet fick köras alls ---
\textcquote[34]{Plass2015Variables}{Students seem to look at it more as a
\enquote{statement of equality}.  If this statement is correct, the
program is executed.  If the statement is not correct, the program
yields an Error}.  Att båda riktningarna är belagda är skälet till att
materialet inte nöjer sig med att visa syntaxfelet, utan säger ut vad de
två operatorerna \emph{gör}.

\paragraph{Exakt en gren, och sedan vidare.}  Att ett falskt villkor
skulle avsluta programmet står i Sorvas inventering av
missuppfattningar hos nybörjare som post nummer 26:
\textcquote[tabell~A.1, s.~360]{Sorva2012Dissertation}{A false condition
ends program if no else branch}.  Den motsatta missuppfattningen --- att
båda grenarna körs, även när villkoret är sant --- står som post nummer
27 i samma tabell:
\textcquote[tabell~A.1, s.~360]{Sorva2012Dissertation}{Both then and else
branches are executed}.  Orden \foreignlanguage{english}{then} och
\foreignlanguage{english}{else} är Pascals; i Python motsvaras
\foreignlanguage{english}{then}-grenen av raderna direkt under
\mintinline{python}{if}.  De två är varandras
spegelbilder och drabbar samma kritiska aspekt: att villkoret
\emph{väljer} gren och varken avslutar programmet eller kör båda.  Därför
är kontrastparet i \cref{ex:varmt-kallt} byggt så att den sista
utskriften bär båda aspekterna på en gång.

\paragraph{Kedjade villkor.}  Att skriva
\mintinline{python}{if x != a or b} i stället för
\mintinline{python}{if x != a or x != b} är dokumenterat i den hittills
största samlingen av nybörjarförklaringar, drygt 16\,000 stycken
insamlade under tre år på en webbplats för programvisualisering.
\Textcite{GuoMarkelZhang2020} tar upp problemet så här:
\textcquote{GuoMarkelZhang2020}{Incorrectly chaining conditions, like
\enquote{if x != a or b} which reads naturally in English as \enquote{if
x is not equal to a or b}}.  Felet kommer alltså ur att villkoret läses
högt.  Samma källa noterar att en slingvariabel ofta
används utan att ha satts före slingan, i tron att den sätts och räknas
upp av sig själv.

\paragraph{Slingor.}  Att en slinga vars styrvariabel aldrig ändras
aldrig tar slut är också ett katalogiserat antimönster:
\textcquote{Bosse2021Antipatterns}{The program will go into an infinite
loop because, since the control variable is never changed, the result of
the condition that controls the repetition will not change either}.
Samma katalog tar upp det omvända valet av konstruktion:
\textcquote{Bosse2021Antipatterns}{Using repetition where only selection
would be required will cause wrong results}.  Det är exakt kontrasten
mellan de två rabattfunktionerna i \cref{ex:rabatt}.  Att
slingor över huvud taget hör till det svåra är belagt i en studie av
nybörjares missuppfattningar om slingor i en grundkurs:
\textcquote{Eckert2022Loops}{For novice programmers loops generally and
nested loops specifically pose big difficulties}.

\ifbool{ltnote}{%
\paragraph{Kontrast före generalisering.}  Att exemplen står parvis, och
att kontrasten kommer före generaliseringen, följer av variationsteorin.
\Textcite[45--47]{Marton2015} visar att en aspekt bara kan urskiljas om
den varierar mot en invariant bakgrund, och att den variationen måste
upplevas samtidigt: \textcquote[46]{Marton2015}{if you cannot see the
greenness of any of three green, though otherwise different, things, you
cannot see what they have in common, either}.
}{}

\section{Diskussion och begränsningar}
\label{sec:cond-diskussion}

Fyra saker begränsar svaret.

För det första är en av källorna läst enbart i sammanfattning, inte i
fulltext: studien om slingor, vars förlag håller texten stängd.  Den
citeras därför bara för att slingor är svåra, vilket sammanfattningen
säger rakt ut, inte för någon detalj i hur svårigheterna mättes.  En
källa som först var påtänkt --- en litteraturöversikt över nybörjares
missuppfattningar --- ströks: det enskilda påstående den skulle bära, att
båda grenarna tros köras, står inte i dess sammanfattning, och
fulltexten gick inte att komma åt.  Påståendet bärs i stället av den
inventering där det står ordagrant.  En källa man inte har läst räknas
inte, hur passande titeln än är.

För det andra säger ingen av källorna något om hur ofta felen
förekommer i just den här kursen.  De säger att felen finns och
återkommer i flera länder och programspråk; de säger inte att just våra
studenter gör dem lika ofta.  Kursens egna diagnostiska frågor finns
till för det.

För det tredje bygger antimönsterkatalogen på kod från en
grundkurs i ett annat land och delvis ett annat språk (C såväl som
Python).  Att felen känns igen här är inte samma sak som att frekvenserna
överförs.

För det fjärde är samstämmighet inte bevis.  Att flera oberoende slags
material --- studenters egna förklaringar, granskad kod, intervjuer och
en inventering --- pekar åt samma håll är starkare än en enda studie,
men ingen av dem är ett experiment, och påståendet materialet gör är
svagare än så: att det \emph{är värt att undervisa mot} de här
missuppfattningarna, inte att de drabbar alla.

\section{Slutsats}
\label{sec:cond-slutsats}

Ja, missuppfattningarna är belagda, med det förbehållet att ett av
påståendena vilar på en källa som ännu bara lästs i sammanfattning, och
att beläggen säger att felen är vanliga utan att säga hur vanliga de är
hos oss.  Att ett likhetstecken skrivs där en jämförelse avses är
dokumenterat i granskad studentkod\autocite{Bosse2021Antipatterns}, och
att en tilldelning läses som ett likhetstest i
intervjuer\autocite{Plass2015Variables}.  Att ett falskt villkor tros
avsluta programmet, och att båda grenarna tros köras, står i samma
inventering av missuppfattningar\autocite[tabell~A.1,
s.~360]{Sorva2012Dissertation}.
Att villkor kedjas utan att jämförelsen upprepas, och att slingvariabeln
glöms bort, är dokumenterat i studenters egna
förklaringar\autocite{GuoMarkelZhang2020}; att en oändlig slinga följer
av en oförändrad styrvariabel och att upprepning väljs där ett val
räcker, i antimönsterkatalogen\autocite{Bosse2021Antipatterns}; och att
slingor är svåra i en studie av en grundkurs, läst i
sammanfattning\autocite{Eckert2022Loops}.%
\ifbool{ltnote}{  Att materialets exempel står
parvis, med kontrasten före generaliseringen, följer av
variationsteorins villkor för
urskiljning\autocite[45--47]{Marton2015}.}{}  Beläggen används till det de
räcker till: att välja \emph{vilka} kontraster som är värda tid på en
föreläsning --- inte till att påstå hur många i salen som har
missuppfattningen.


\section{Fortsatt arbete}
\label{sec:cond-fortsatt}

Det största som saknas är en egen sökning.  Bilagan vilar på källor ur
en genomgång gjord för en annan fråga (\cref{sec:cond-metod}), och
därför finns här varken sökfrågor, motsökning för just de här
formuleringarna eller en träfflista att granska; att ställa den
tvåsidiga frågan för vart och ett av de \ifbool{ltnote}{fem}{fyra}
påståendena, och redovisa
den med tabell och träfflista, är nästa steg.  Två källor är dessutom
ofärdiga.  Studien om slingor är läst enbart i
sammanfattning\autocite{Eckert2022Loops}, eftersom förlaget håller
texten stängd; fulltexten skulle säga vilka missuppfattningar om
slingor som faktiskt hittades, och källan skulle då kunna bära mer än
att slingor är svåra.  Och den översikt över nybörjares
missuppfattningar som skulle ha burit påståendet om att båda grenarna
körs --- Qian och Lehman, 2017, i \foreignlanguage{english}{ACM
Transactions on Computing Education}\footnote{\textsc{doi}
10.1145/3077618.} --- gick inte att komma åt och är därför inte använd
(\cref{sec:cond-diskussion}); en läsning av den skulle visa om
påståendet vilar på mer än den inventering som nu bär det.

För ett säkrare svar saknas mätningar i den population materialet
vänder sig till.  Missuppfattningarna är dokumenterade i studentkod,
intervjuer och inventeringar från andra kurser, andra länder och delvis
andra programspråk; ingen av källorna säger hur vanliga de är här.  En
räkning i kursens egna inlämningar och svar skulle säga det.  Otestat
är också om missuppfattningarna hänger ihop --- om samma student
som läser en tilldelning som ett likhetstest också tror att ett falskt
villkor avslutar programmet --- vilket avgör om ett enda kontrastpar
kan träffa flera av dem eller om var och en behöver sitt.  Visar en
sådan undersökning att felen är sällsynta här, kan materialet lägga
tiden på färre kontraster; visar den att missuppfattningarna hänger
ihop, kan kontrastparen slås samman till färre och tydligare.

\ifbool{ltnote}{%
Ett påstående i materialet är inte prövat alls.  Programmen skrivs
genomgående som funktioner, och skälet till regeln, som lärarnoterna i
föreläsningen \emph{Funktioner} anger, är att ett material som inte
följer sin egen rekommendation får den att framstå som frivillig.  Det
är ett empiriskt påstående om hur studenter uppfattar ett
kursmaterial, och ingen sökning har gjorts för det; den tvåsidiga
frågan --- lär sig studenter en praxis bättre när exemplen genomgående
följer den, och finns det studier som visar motsatsen? --- är en egen
bilaga som återstår att skriva.
}{}

% Author's note (verification and reproduction):
%   sources           : reused from the author's misconceptions study
%                       (~/devel/edu/research/vt-prog-misconceptions,
%                       conditionals.tex and repetitions.tex; literature
%                       protocol app:litprotocol), cited as SooriBosk2026
%                       (entry copied from modules/variables/slides;
%                       2026-09-16, replacing the nickname "den egna
%                       artikeln om missuppfattningar" in sec:cond-metod
%                       and in four ltnotes), plus Marton2015 and
%                       HuntThomas1999 reused from
%                       modules/computational-thinking/slides/ltnotes.bib
%   no new search     : deliberately none; see sec:cond-metod.  If a new
%                       claim is added to this deck, run the two-sided
%                       protocol under a session named
%                       conditionals-<claim-slug> and add a chapter with a
%                       query table and a classified hit list.
%   provenance        : ltnotes.bib, keys Bosse2021Antipatterns
%                       Plass2015Variables Sorva2012Dissertation
%                       MisconceptionsSurvey2017 GuoMarkelZhang2020
%                       Eckert2022Loops Marton2015 HuntThomas1999
%                       Martin2003Agile SooriBosk2026
%                       (CLAIM / FOUND-VIA / PICKED / QUOTE / VERIFIED /
%                       DATE, with "(here)" lines for this deck)
%   abstract-level    : Eckert2022Loops (IEEE Xplore closed)
%   dropped            : MisconceptionsSurvey2017 (Qian & Lehman 2017).
%                       It was to carry the both-branches item, which is a
%                       cross-cite in the misconceptions paper and is NOT
%                       in the ACM abstract; the ACM full text is
%                       Cloudflare-gated (the download returns an HTML
%                       error page).  Replaced by Sorva Table A.1 no. 27,
%                       verified verbatim in the aaltodoc full text
%                       (printed p. 360 = PDF page 368) on 2026-09-03.
%   already backed    : DRY in modules/computational-thinking/slides,
%                       bilaga E; SRP in modules/variables/
%                       slides-functions, bilaga B (KISS is no longer
%                       mentioned in this deck, 2026-09-22).
